% Part: lambda-calculus
% Chapter: lambda-definability
% Section: fixpoints

\documentclass[../../../include/open-logic-section]{subfiles}

\begin{document}

\olfileid{lam}{ldf}{fp}

\olsection{不动点}

假设我们想通过递归将阶乘函数定义为一个满足以下性质的项 $\fn{Fac}$：
\[
\fn{Fac} \ident \lambd[n][\fn{IsZero}\, n\, \num 1 (\fn{Mult}\, n (\fn{Fac}(\fn{Pred} \, n)))]
\]
也就是说，如果 $n = 0$，则 $n$ 的阶乘是 $1$；否则就是 $n$ 乘以 $n-1$ 的阶乘。
当然，我们不能这样定义项 $\fn{Fac}$，因为 $\fn{Fac}$ 本身出现在了等式右边。
这种涉及自引用的递归定义，并不属于λ演算的一部分。
例如，通过
\[
\fn{Mult} \ident \lambd[ab][a (\fn{Add}\, a) 0]
\]
来定义一个项时，等式右边只涉及先前定义的项，如 $\fn{Add}$。
我们总可以通过将 $\fn{Add}$ 替换为定义它的项来消除它。
这样就能将 $\fn{Mult}$ 表示为一个纯λ项；如果 $\fn{Add}$ 本身也包含已定义的项（如 $\fn{Add}'$ 那样），我们可以继续这个过程，最终得到一个纯λ项。

然而，对于上面 $\fn{Fac}$ 这样的递归定义，情况并非如此。
如果我们将右边出现的 $\fn{Fac}$ 替换为 $\fn{Fac}$ 本身的定义，会得到：
\begin{align*}
  \fn{Fac} & \ident \lambd[n][\fn{IsZero}\, n\, \num{1}] \\
    & \qquad (\fn{Mult}\, n
  ((\lambd[n][\fn{IsZero} \, n \, \num 1\, (\fn{Mult}\, n\, (\fn{Fac}
    (\fn{Pred}\, n)))]) (\fn{Pred}\, n)))
\end{align*}
而我们依然没有消除掉右边的 $\fn{Fac}$。
显然，如果我们重复这个过程，定义会变得越来越长，且该过程永远无法产生一个纯λ项。
因此，这种定义阶乘（或者更一般的递归函数）的方法是行不通的。

不过，这个递归定义确实告诉了我们一些信息：如果 $f$ 是表示阶乘函数的项，那么将项
\[
\fn{Fac}' \ident \lambd[g][\lambd[n][\fn{IsZero} \, n \, \num 1\, (\fn{Mult} \, n\, (g (\fn{Pred} n)))]]
\]
应用于项 $f$（即 $\fn{Fac}'\,f$），也表示阶乘函数。
也就是说，如果我们把 $\fn{Fac}'$ 看作一个以函数为自变元并返回函数的函数，那么只要 $f$ 是阶乘函数，$\fn{Fac}'\, f$ 的值就是~$f$。
一个满足 $\fn{Fac}' \, f \equal[\beta] f$ 的函数~$f$，称为 $\fn{Fac}'$ 的一个\emph{不动点}。所以，阶乘函数就是 $\fn{Fac}'$ 的一个不动点。

λ演算中存在一些项，它们能计算给定项的不动点，这些项可以用来由 $\fn{Fac}'$ 这样的项定义阶乘。

\begin{defn}
  \ollabel{defn:Turing-Y}
  \emph{Y组合子}是如下项：
  \[
  Y \ident (\lambd[ux][x(uux)])(\lambd[ux][x(uux)]).
  \]
\end{defn}

\begin{thm}
  对于任何项 $g$，$Y$ 都具有性质 $Yg \red g(Yg)$。因此，$Yg$ 总是~$g$ 的一个不动点。
\end{thm}

\begin{proof}
  用~$U$ 缩写 $(\lambd[ux][x(uux)])$，于是 $Y \ident UU$。那么
  \begin{align*}
    Y g &\ident (\lambd[ux][x(uux)])U\,g \\
    &\red (\lambd[x][x(UUx)])g\\
    & \red g(UUg) \ident g(Yg).
  \end{align*}
  由于 $g(Yg)$ 和 $Yg$ 都可归约为 $g(Yg)$，故 $g(Yg) \equal[\beta] Yg$，所以 $Yg$ 是~$g$ 的一个不动点。
\end{proof}

当然，由于 $Yg$ 是一个可约式，归约可以无限进行下去：
\begin{align*}
  Y g &\red g (Y g) \\
    &\red g (g (Y g)) \\
    &\red g(g (g (Y g)))\\
    &\ldots
\end{align*}
所以我们可以把 $Yg$ 看作是 $g$ 无限次应用于自身的结果。
如果我们再对它额外应用一次~$g$，打个比方说，并没有做任何额外的事情；将~$g$ 作用于已对~$Yg$ 无限次作用的~$g$，结果仍是对~$Yg$ 无限次作用~$g$。

注意，从~$Yg$ 开始的上述β归约步骤序列是无限的。
因此，如果我们把 $Yg$ 应用于某个项，即考虑 $(Yg)N$，那么这个项也会归约到无限多个不同的项，即 $(g(Yg))N$, $(g(g(Yg)))N$, ……。
尽管如此，某些\emph{其他}的归约步骤序列还是有可能终止于一个范式。

以阶乘为例。将 $\fn{Fac}$ 定义为 $Y\, \fn{Fac}'$（即 $\fn{Fac}'$ 的一个不动点）。那么：
\begin{align*}
  \fn{Fac}\,\num 3
  &\red Y\, \fn{Fac}' \, \num 3 \\
  &\red \fn{Fac}' (Y \,\fn{Fac}') \, \num 3 \\
  &\ident (\lambd[x][\lambd[n][\fn{IsZero} \, n\, \num 1\,
      (\fn{Mult}\, n\, (x (\fn{Pred}\, n)))]]) \, \fn{Fac} \, \num 3\\
  & \red \fn{IsZero} \, \num 3\, \num 1\,
      (\fn{Mult}\, \num 3\, (\fn{Fac} (\fn{Pred}\, \num 3))) \\
  &\red  \fn{Mult}\, \num 3 \, (\fn{Fac} \, \num 2).
  \intertext{类似地，}
  \fn{Fac}\,\num 2
  &\red  \fn{Mult}\, \num 2 \, (\fn{Fac} \, \num 1) \\
  \fn{Fac}\,\num 1
  &\red  \fn{Mult}\, \num 1 \, (\fn{Fac} \, \num 0)
  \intertext{但是}
  \fn{Fac}\,\num 0
  &\red \fn{Fac}' (Y \,\fn{Fac}') \, \num 0 \\
  &\ident (\lambd[x][\lambd[n][\fn{IsZero} \, n\, \num 1\,
      (\fn{Mult}\, n\, (x (\fn{Pred}\, n)))]]) \, \fn{Fac} \, \num 0\\
  & \red \fn{IsZero} \, \num 0\, \num 1\,
      (\fn{Mult}\, \num 0\, (\fn{Fac} (\fn{Pred}\, \num 0))).\\
  &\red \num 1.
  \intertext{因此合起来}
  \fn{Fac}\,\num 3
  &\red  \fn{Mult}\, \num 3 \,
  (\fn{Mult} \, \num 2\, (\fn{Mult} \, \num 1\, \num 1)).
\end{align*}

适用于 $\fn{Fac}'$ 的情况同样适用于任何递归定义。
假设我们有一个递归方程
\begin{align*}
g\,x_1\dots x_n & \equal[\beta] N
\intertext{其中 $N$ 可能包含~$g$ 和 $x_1$, \dots,~$x_n$。那么总存在一个项~$G \ident (Y \lambd[g][\lambd[x_1\dots x_n][N]])$，使得}
G\,x_1 \dots x_n & \equal[\beta] \Subst{N}{G}{g}.
\intertext{因为根据不动点定理，}
G \ident (Y \lambd[g][\lambd[x_1\dots x_n][N]]) & \red \lambd[g][\lambd[x_1\dots x_n][N]](Y \lambd[g][\lambd[x_1\dots x_n][N]])\\
& \ident (\lambd[g][\lambd[x_1\dots x_n][N]])\,G
\intertext{从而}
G\,x_1 \dots x_n 
& \red (\lambd[g][\lambd[x_1\dots x_n][N]])\,G\,x_1\dots x_n\\
& \red (\lambd[x_1\dots x_n][\Subst{N}{G}{g}])\,x_1\dots x_n\\
  & \red \Subst{N}{G}{g}.
\end{align*}

\olref{defn:Turing-Y} 中的 $Y$ 组合子归功于 Alan Turing。
Alonzo Church 提出了一个不同的版本，我们称之为~$Y_C$：
\[
Y_C \ident \lambd[g][(\lambd[x][g(xx)])(\lambd[x][g(xx)])].
\]
Church 的组合子比 Turing 的稍弱一些，因为 $Yg \equal[\beta] g(Yg)$ 但并非 $Yg \bred g(Yg)$。
令 $V$ 为项 $\lambd[x][g(xx)]$，使得 $Y_C \ident \lambd[g][VV]$。那么
\begin{align*}
  VV & \ident (\lambd[x][g(xx)])V \red g(VV)
  \text{ 因此}\\
  Y_C g & \ident (\lambd[g][VV])g \red VV \red g(VV), \text{ 但也有}\\
  g(Y_C g) & \ident g((\lambd[g][VV])g) \red g(VV).
\end{align*}
换言之，$Y_Cg$ 和 $g(Y_Cg)$ 都归约到同一个项 $g(VV)$；因此 $Y_Cg
\equal[\beta] g(Y_Cg)$。这对于许多应用来说已经足够了。

\end{document}